\section{Payload Format}
\subsection{Format Selection}
Telemetry payloads are serialised as UTF-8 JSON. This prioritises developer ergonomics and debuggability: payloads are human-readable, trivially inspected with any MQTT client, and require no schema negotiation between publisher and subscriber. A compact binary format is used only for the optional full-rate stream (Section~\ref{sec:daq2}).

\subsection{Burst Payload}
The ADC sample rate was fixed at approximately 30~kHz per channel, which resolves the original choice between single-sample and burst publishing in favour of bursts. Every 5~s, each node publishes one 50~Hz cycle of each channel, decimated to approximately 1~kHz: 20 samples per message, one message per channel.

\begin{lstlisting}[language=json]
  {
    "ts": 1791345068887514368,
    "sync": 1,
    "dt_ns": 998424,
    "v": [231.45, 229.80, 224.12, ...]
  }
\end{lstlisting}

\begin{table}[H]
  \caption{Burst Payload Description}\label{tab:burst-payload}
  \centering
  \begin{tabular}{llp{0.6\linewidth}}
    \hline
    \textbf{Field} & \textbf{Type} & \textbf{Description} \\
    \hline
    \texttt{ts} & integer & Unix (UTC) time in nanoseconds of the first sample, from the node's PTP-disciplined clock \\
    \texttt{sync} & integer & 1 if the node's PTP servo was locked when the burst was captured, otherwise 0 \\
    \texttt{dt\_ns} & integer & Interval between samples in nanoseconds \\
    \texttt{v} & number[20] & Sample values in SI units (V or A), to two decimal places \\
    \hline
    \end{tabular}
\end{table}

Sample $k$ was taken at $t_s + k \cdot \Delta t$, where $t_s$ is \texttt{ts} and $\Delta t$ is \texttt{dt\_ns}. The interval is the nearest whole number of ADC frames to 1~ms (for example 998\,424~ns at 30\,048 frames/s), so it is reported explicitly rather than assumed.

Implementation notes:
\begin{itemize}[noitemsep]
  \item The PTP clock runs on TAI. \texttt{ts} is converted to UTC by subtracting the TAI$-$UTC offset (\texttt{currentUtcOffset}) from the grandmaster's Announce messages, or 37~s until one has arrived. No offset is subtracted if the grandmaster announces an arbitrary (ARB) timescale.
  \item \texttt{ts} is the centre of the first sample's accumulated ADC burst, derived from the hardware timer and PTP clock rather than from interrupt timing.
  \item Currents are sampled one ADC slot after the voltages, so each current message's \texttt{ts} is offset from the voltage messages by that interval.
  \item \texttt{sync} is sent as an integer rather than a JSON boolean, because the historian bridge's JSON parser drops booleans.
  \item A burst is started only on data whose timestamp is valid, and is restarted if a frame gap or PTP clock step occurs during collection, so the 20 samples are always contiguous and on a single time base.
\end{itemize}

Units are implied by the topic rather than included in the payload: a message on \path{microgrid/node_01/voltage/phase_a} carries volts.

The single-sample payload of the original design (\texttt{\{"ts", "v"\}}) is not used.

\subsection{PTP Status Payload}
\label{sec:ptpstatus}
On every PTP Sync, about once per second, each node publishes the result of its servo update on \texttt{microgrid/<node\_id>/ptp\_status/servo}:

\begin{lstlisting}[language=json]
  {
    "ts": 1791696297692880523,
    "offset_ns": 18274,
    "path_delay_ns": 15497,
    "freq_ppb": 18010,
    "state": 2,
    "locked": 0,
    "stepped": 0,
    "utc_offset_s": 37,
    "utc_announced": 1
  }
\end{lstlisting}

\begin{table}[H]
  \caption{PTP Status Payload Description}
  \centering
  \begin{tabular}{llp{0.6\linewidth}}
    \hline
    \textbf{Field} & \textbf{Type} & \textbf{Description} \\
    \hline
    \texttt{ts} & integer & The grandmaster's time for that Sync, as Unix (UTC) nanoseconds, using the same TAI offset as the bursts \\
    \texttt{offset\_ns} & integer & Measured clock offset from the grandmaster \\
    \texttt{path\_delay\_ns} & integer & Mean path delay used in the offset calculation \\
    \texttt{freq\_ppb} & integer & Frequency correction applied by the servo, relative to the nominal clock rate \\
    \texttt{state} & integer & Servo state: 0 no Sync yet, 1 estimating frequency, 2 PI loop running \\
    \texttt{locked} & integer & 1 once eight consecutive offsets are within 1~\textmu s \\
    \texttt{stepped} & integer & 1 if this update stepped the clock (offset above 1~ms) \\
    \texttt{utc\_offset\_s} & integer & TAI$-$UTC offset subtracted from \texttt{ts} and the bursts' timestamps \\
    \texttt{utc\_announced} & integer & 1 if that offset came from the grandmaster's Announce, 0 if it is the 37~s default \\
    \hline
  \end{tabular}
\end{table}

The values are 64-bit integers formatted by the firmware itself, since its \texttt{printf} lacks 64-bit support; offsets before the first clock step can exceed $10^{18}$~ns. As for the bursts, booleans are sent as integers.

\subsection{Full-Rate Binary Payload}
\label{sec:daq2}
When built with \texttt{MQTT\_PUBLISH\_DAQ2} (off by default), each node additionally publishes every acquisition packet on \texttt{daq2/<node\_id>/samples}. The payload is a 128-byte little-endian header followed by \textit{frames} $\times$ \textit{channels} \texttt{float32} values in row order. The header carries a sequence number and frame index (so gaps reveal dropped data), the timestamp of the first frame and the frame period, per-channel sampling offsets, the PTP epoch (which changes on every clock step), a timing uncertainty estimate, and flags indicating PTP lock and timestamp validity. The flags also carry the grandmaster's timescale and, in their upper 16 bits, the TAI$-$UTC offset the node received in its Announce messages. The full layout is defined by \texttt{daq\_packet\_t} in the firmware's \texttt{app/daq/daq\_task.h}. Publishing is best-effort: a packet that cannot be queued is dropped, as sampling never waits on the network. The stream is ingested by the \texttt{daq2-ingest} service, described in the InfluxDB design document.
